% !TeX root = ../Cosmology.tex

%% --------------------------------------------------------
\chapter{Сучасні космологічні дані та параметри раннього Всесвіту}
\label{app:data}
%% --------------------------------------------------------

У цьому додатку зведено ключові космологічні параметри, отримані на основі комбінації сучасних спостережних даних: кутового спектра анізотропії реліктового випромінювання (CMB) космічної обсерваторії \textit{Planck} (реліз 2018 року / Legacy data), даних баріонних акустичних осциляцій (BAO) з оглядів галактик (SDSS/eBOSS)~\cite{eisenstein_2005}, а також найсвіжіших локальних вимірювань та обмежень раннього Всесвіту.

\section{Ключові параметри стандартної космологічної\\ моделі ($\Lambda$CDM)}

У таблиці~\ref{tab:cosmo_parameters} наведено базові (6-параметричні) та похідні космологічні характеристики Всесвіту на сьогоднішній день для комбінації даних \textit{Planck} TT,TE,EE + lowE + lensing + BAO (довірчий інтервал $68\%$)~\cite{planck_2018}.

\begin{table}[h!]
\centering
\caption{Сучасні значення космологічних параметрів $\Lambda$CDM~\cite{planck_2018}.}
\label{tab:cosmo_parameters}
\begin{tblr}{
colspec={Q[l, m, 7cm]cc}
}
\hline
\textbf{Параметр} & \textbf{Позначення} & \textbf{Значення (68\% CL)} \\ \hline
\SetCell[r=1, c=3]{c,m}{\textit{Базові (незалежні) параметри:}} \\
Густина баріонів & $\Omega_b h^2$ & $0.02242 \pm 0.00014$ \\
Густина холодної темної матерії & $\Omega_c h^2$ & $0.11933 \pm 0.00091$ \\
Кутовий масштаб звукового горизонту & $100\,\theta_{\rm MC}$ & $1.04109 \pm 0.00041$ \\
Оптична товща реонізації & $\tau$ & $0.0544 \pm 0.0073$ \\
Спiральний індекс скалярних збурень & $n_s$ & $0.9665 \pm 0.0038$ \\
Амплітуда первинного спектра кривини & $\ln(10^{10} A_s)$ & $3.044 \pm 0.014$ \\ \hline
\SetCell[r=1, c=3]{c,m}{\textit{Похідні параметри:}} \\
Сучасне значення сталої Хаббла & $H_0$ & $67.4 \pm 0.5$ км/с/Мпк \\
Повний параметр густини матерії & $\Omega_m$ & $0.315 \pm 0.007$ \\
Параметр густини темної енергії ($\Lambda$) & $\Omega_\Lambda$ & $0.685 \pm 0.007$ \\
Флуктуації густини на масштабі $8\,h^{-1}$ Мпк & $\sigma_8$ & $0.811 \pm 0.006$ \\
Вік Всесвіту & $t_0$ & $13.787 \pm 0.020$ млрд років \\ \hline
\end{tblr}
\end{table}

\textit{Примітка щодо напруги Хаббла (Hubble Tension):} Наведене значення $H_0 \approx 67.4$ км/с/Мпк є результатом моделювання «зверху-вниз» на основі раннього Всесвіту (CMB)~\cite{planck_2018}одом «космічної драбини відстаней» проекту SH0ES (на основі Цефеїд та Надновах типу Ia), які дають $H_0 = 73.04 \pm 1.04$ км/с/Мпк.

\section{Параметри та обмеження на стадію інфляції}

Первинні збурення, зафіксовані в CMB, накладають жорсткі обмеження на енергетичний масштаб інфляції та форму потенціалу інфлатона $V(\phi)$. Динаміку інфлатона і параметри повільного кочення розглянуто у розд.~\ref{sec:inflaton}; первинний спектр збурень і визначення спектрального індексу~--- у розд.~\ref{sec:quantum}.
\begin{itemize}
    \item \textbf{Тензорно-скалярне відношення ($r$):} Сумісні дані \textit{Planck 2018}~\cite{planck_2018} та масивів BICEP/Keck~\cite{bicep_keck_2021} (обмеження за B-модою поляризації, розд.~\ref{sec:pol}) встановлюють верхню межу:
    \[ r_{0.002} < 0.032 \quad (95\% \text{ CL}) \]
    Відношення $r = 16\varepsilon$ виражається через параметр повільного кочення $\varepsilon$~\eqref{eq:eps} за формулою~\eqref{eq:tensor_scalar_ratio}. Це повністю виключає прості степеневі потенціали $V(\phi) \propto \phi^2$ (хаотична інфляція Лінде, розд.~\ref{sec:chaotic}) та робить фаворитами моделі типу Старобінського ($R^2$-гравітація)~\cite{starobinsky_1980} або інфляцію Гіггса~\cite{bezrukov_2008}; пор.~також табл.~\ref{tab:inflation_models}.

    \item \textbf{Енергетичний масштаб інфляції ($V_{\rm inf}$):} Верхня межа на $r$ обмежує густину енергії під час повільного кочення. З формули для тензорного спектра~\eqref{eq:tensor_power_spec} та умови~\eqref{eq:sr_approximations} ($H^2 \approx V/3M_{\rm Pl}^2$) випливає:
    \[ V^{1/4} < 1.6 \times 10^{16} \text{ GeV} \quad (95\% \text{ CL}) \]
    Це вказує на те, що інфляція відбувалася поблизу масштабу Великого Об'єднання (GUT scale).

    \item \textbf{Спектральний індекс ($n_s$):} Значення $n_s = 0.9665 \pm 0.0038$ (табл.~\ref{tab:cosmo_parameters}) пов'язане з параметрами slow-roll формулою~\eqref{eq:ns}: $n_s - 1 = -2\varepsilon - \eta_V$, де $\varepsilon$~\eqref{eq:eps} та $\eta_V$~\eqref{eq:slow_roll_pot} обчислюються в момент виходу моди за горизонт.

    \item \textbf{Просторова кривина Всесвіту ($\Omega_K$):} Комбінація спектра CMB з ефектом гравітаційного лінзування та BAO~\cite{eisenstein_2005} підтверджує плоску геометрію (рівняння Фрідмана~\eqref{eq:F1}--\eqref{eq:F2} та $\Omega_K$ у розд.~\ref{sec:friedmann}):
    \[ \Omega_K = 0.0007 \pm 0.0019 \]
    Це повністю узгоджується з інфляційним передбаченням тривалості стадії $N_{\rm total} \gtrsim 50-60$, яка виводиться у додатку~\ref{sec:appendix_efolds} (формула~\eqref{eq:app_final_n}).
\end{itemize}

\section{Енергетичні та температурні характеристики стадії розігріву (Preheating / Reheating)}

Оскільки стадія пре-розігріву є сильно нелінійним квантовим процесом, безпосередньо виміряти її температуру через релікти важко. Проте сучасна космологія накладає суворі \textit{теоретичні та феноменологічні межі} на основі узгодження первинного нуклеосинтезу (BBN) та тривалості інфляції (параметра $N_{\rm CMB}$). Детальний розгляд механізму розігріву та параметричного резонансу наведено у розд.~\ref{sec:reheating}~\cite{kofman_1994,kofman_1997}:

\begin{itemize}
    \item \textbf{Абсолютна нижня межа температури розігріву ($T_{\rm rh}^{\rm min}$):}
    \[ T_{\rm rh} \gtrsim 4-5 \text{ MeV} \]
    Ця межа є залізобетонною. Всесвіт \textit{мусив} стати радіаційно-до\-мі\-но\-ва\-ним і термалізованим до моменту, коли його вік становив $\sim 1$ секунду. Якщо $T_{\rm rh}$ була б нижчою, параметри розпаду нейтронів та швидкість розширення порушили б первинний нуклеосинтез, що суперечило б спостережуваній кількості легких елементів ($^4\text{He}$, $\text{D}$, $^7\text{Li}$).

    \item \textbf{Верхня межа температури на момент Preheating ($T_{\rm max}$):}
    Енергія коливань інфлатона відразу після закінчення повільного кочення задає максимальну можливу температуру. З першого рівняння Фрідмана~\eqref{eq:F1} та виразу для густини~\eqref{eq:inflaton_rho_p} за умови $V_{\rm end} \approx \rho_{\rm inf}$~\cite{kofman_1997}:
    \[ T_{\rm max} \sim \left( \frac{30}{\pi^2 g_*} V_{\rm end} \right)^{1/4} \lesssim 10^{15} - 10^{16} \text{ GeV} \]
    де $g_* \sim 100-200$ --- ефективна кількість релятивістських ступенів вільності (пор.~\eqref{eq:app_hinf}).

    \item \textbf{Ефективна температура на стадії параметричного резонансу\\ (пре-розігрів):}
    Під час вибухового народження часток через рівняння Матьє~\eqref{eq:mathieu} (розд.~\ref{sec:reheating}) система перебуває у \textit{нетермальному стані}~\cite{kofman_1994}. Поняття «температури» в термодинамічному сенсі ще не існує, оскільки спектри часток $\chi_k$ є високоселективними та анізотропними. Проте, якщо оцінити еквівалент густоти енергії $\rho_\chi \sim g^2 \Phi^2 \chi^2$, квазі-тем\-пе\-ра\-тур\-ний масштаб первинного спалаху становить:
    \[ T_{\rm inst} \sim 10^{13} - 10^{14} \text{ GeV} \]

    \item \textbf{Температура завершення розігріву ($T_{\rm rh}$):}
    Повна термалізація настає при $H \approx \Gamma_\phi$. Підставляючи $\rho_r = \frac{\pi^2}{30}g_* T^4$ у перше рівняння Фрідмана~\eqref{eq:F1}, отримуємо формулу~\eqref{eq:t_reheating}:
    \[ T_{\rm rh} = \left(\frac{90}{8\pi^3 g_*}\right)^{\!1/4} \sqrt{M_{\rm Pl}\,\Gamma_\phi} \]

    \item \textbf{Обмеження з проблеми Гравітіно (Gravitino bound):}
    У суперсиметричних теоріях (SUSY) занадто висока температура $T_{\rm rh}$ призводить до надлишкового народження стабільних або метастабільних гравітіно, які руйнують легкі ядра під час BBN~\cite{kofman_1997}. Щоб цього уникати, накладається верхня межа:
    \[ T_{\rm rh} \lesssim 10^9 - 10^{10} \text{ GeV} \]
\end{itemize}

\section{Реліктовий фон та інші компоненти на сьогодні}

\begin{itemize}
    \item \textbf{Температура CMB на сьогодні ($T_{\rm CMB}$):} Виміряна з найвищою точністю приладом FIRAS на супутнику COBE і уточнена пізнішими місіями~\cite{planck_2018}. Вона є початковою умовою у виразі~\eqref{eq:app_entropy} для $T_0$:
    \[ T_{\rm CMB, 0} = 2.72548 \pm 0.00057 \text{ K} \]
    \item \textbf{Густина реліктових фотонів:} $n_\gamma \approx 411 \text{ см}^{-3}$.
    \item \textbf{Температура реліктових нейтрино (передбачувана):} \\
    $T_\nu = (4/11)^{1/3} T_{\rm CMB} \approx 1.945 \text{ K}$, що отримується з закону збереження ентропії~\eqref{eq:app_entropy}; внесок нейтрино у $H(z)$ пропорційний $\Omega_\nu$.
    \item \textbf{Обмеження на суму мас нейтрино ($\sum m_\nu$):} Комбінація \textit{Planck} CMB + lensing + BAO~\cite{planck_2018,eisenstein_2005} обмежує сумарну масу трьох активних типів нейтрино зверху:
    \[ \sum m_\nu < 0.12 \text{ eV} \quad (95\% \text{ CL}) \]
    Це вказує на те, що внесок нейтрино у сучасну густину речовини є мізерним ($\Omega_\nu \lesssim 0.003$); критична густина $\rho_{\rm crit}$ визначена у~\eqref{eq:rho_crit_new}.
\end{itemize}